\title{FlashAttention-3: Fast and Accurate Attention with Asynchrony and Low-precision}

\begin{abstract}
As a fundamental element in the widely used Transformer framework, attention represents a primary computational constraint in scaling up large language models and managing extended context windows. The \textsc{FlashAttention} algorithm enhanced attention efficiency on GPUs by strategically reducing memory operations. Nonetheless, existing approaches have not fully utilized the advanced features of contemporary hardware; notably, \textsc{FlashAttention-2} attains just 35\% hardware utilization on the H100 GPU. In this work, we introduce three principal strategies to accelerate attention mechanisms on Hopper GPUs: leveraging Tensor Core and TMA asynchrony to (1) concurrently schedule computation and data transfer through warp specialization, (2) interleave block-level matrix multiplications with softmax calculations, and (3) incorporate block-wise quantization and non-sequential processing, supported by hardware for FP8 low-precision arithmetic. Our approach, \textsc{FlashAttention-3}, realizes a speed increase of 1.5-2.0$\times$ on H100 GPUs, achieving up to 740 TFLOPs/s (75\% of peak throughput) in FP16 mode and approaching 1.2 PFLOPs/s with FP8 precision. Experimental results indicate that FP8 \textsc{FlashAttention-3} reduces numerical errors by a factor of 2.6$\times$ compared to a reference FP8 attention method.
\end{abstract}

\section{Introduction}
\label{sec:intro}

Within the Transformer architecture~\citep{vaswani2017attention}, the primary limitation in computation arises from the attention mechanism, as determining the self-attention scores between queries and keys grows quadratically with respect to sequence length. Enabling attention over more extended contexts stands to facilitate capabilities such as processing and reasoning across multiple lengthy documents~\citep{guo2021longt5,shaham2022scrolls,peng2023yarn} and analyzing large codebases containing numerous files~\citep{roziere2023code, li2023starcoder}. It further extends to support additional modalities, including high-resolution imagery~\citep{chen2022scaling}, audio~\citep{gulati2020conformer}, and video data~\citep{ho2022video}, as well as innovative applications like user interaction tracking with extensive historical data~\citep{sun2019bert4rec} and agent planning across long-term horizons~\citep{yao2022react}. Consequently, there has been growing momentum toward accelerating attention in long-context scenarios—this includes techniques such as approximation~\citep{katharopoulos2020transformers,choromanski2020rethinking, tay2020efficient}, software-level improvements (\citep{rabe2021self,dao2022flashattention,kwon2023efficient}), and the exploration of alternative model structures~\citep{peng2023rwkv,sun2023retentive,gu2023mamba}.

This study extends the foundational contributions of \citet{dao2022flashattention}, who developed exact-attention methods informed by the GPU's execution framework and hardware intricacies. In their publication, Dao et al. proposed \textsc{FlashAttention}, featuring an innovative tiling mechanism for efficient parallelization of attention. This approach consolidates all attention-related operations into a unified GPU kernel, thereby circumventing costly intermediate transfers to global memory. 

Building upon this, \citet{dao2023flashattention2} transformed the algorithm into \textsc{FlashAttention-2}, adding the capability to parallelize computations across the sequence length axis. They accomplished this by segmenting the forward pass’s core loop across blocks of the key and value matrices, which consequently enhances GPU occupancy and workload distribution. Despite these advancements, our assessment reveals that \textsc{FlashAttention-2} still demonstrates suboptimal resource utilization on the latest GPU models when compared to highly tuned matrix-multiplication (GEMM) kernels—achieving, for instance, only about 35\% usage versus the 80-90\% attained by GEMM kernels on the Hopper H100. This performance discrepancy is partly due to implementation-specific factors, such as the reliance on Ampere instructions for Tensor Core operations instead of exploiting Hopper-optimized instructions. Prior studies like ThunkerKitten~\citep{spector2024thunder} and cuDNN 9~\citep{cudnn9} have indicated that leveraging Hopper-centric instructions in conjunction with tiled abstractions can accelerate attention processing and lead to simpler, more efficient implementations.

At a foundational level, the algorithm underpinning \textsc{FlashAttention-2} operates within a straightforward synchronous framework and deliberately omits considerations for asynchronous execution or low-precision arithmetic in its design. Asynchronous processing emerges due to specialized hardware components optimized for key tasks in machine learning workflows: certain units focus on matrix multiplication (Tensor Cores) or memory access (Tensor Memory Accelerator -- TMA), while others, such as CUDA cores, handle logical, integer, and floating-point operations. The adoption of reduced precision formats—including FP8 with Hopper and FP4 with Blackwell, following earlier usage of FP16 (Pascal, 2017) and BF16 (Ampere, 2020)—has reliably resulted in two- or four-fold improvements in throughput without incurring additional costs in power consumption or silicon footprint. Section~\cref{subsec:hardware} provides a survey of these advancements in Hopper. The main engineering problem involves rearchitecting \textsc{FlashAttention-2} to take advantage of these specialized hardware functionalities: harnessing asynchrony necessitates concurrent execution of matmul and softmax despite inherent output dependencies, while employing low-precision computation requires diligent strategies to control quantization errors, especially for outlier features characteristic of LLMs~\citep{dettmers2208llm, sun2024massive}.

In pursuit of this objective, we introduce \textsc{FlashAttention-3}, which integrates three innovative approaches to enhance efficiency on modern GPU platforms:\footnote{While our findings are presented with reference to NVIDIA's Hopper architecture, the underlying algorithm remains applicable to any GPU architecture that supports strong asynchronous execution and offers low-precision features.}

\begin{enumerate}

\item \textbf{Producer-Consumer asynchrony:} We introduce a warp-level software pipelining strategy that leverages asynchronous operations by assigning separate warps to the producer and consumer tasks involved in data transfer and Tensor Core computation. This separation enhances the algorithm’s capacity to mask both memory access delays and instruction dispatch latencies.

\item \textbf{Masking softmax within asynchronous block-wise GEMMs:} We enable the overlap of low-throughput non-GEMM computations required in softmax (such as floating point multiply-add and exponential functions) with the execution of WGMMA instructions for GEMM. To achieve this, we modify the \textsc{FlashAttention-2} algorithm to eliminate specific serial dependencies between the softmax steps and GEMM operations. For instance, in the algorithm’s two-stage variant, softmax is computed on a score matrix block simultaneously as WGMMA asynchronously processes the next block.

\item \textbf{Low-precision GEMM with hardware acceleration:} We reconfigure the forward pass to utilize FP8 Tensor Cores for GEMM, substantially increasing observed TFLOPs/s by nearly twofold. This adaptation requires reconciling the distinct memory layout expectations for blocks of FP32 accumulators versus FP8 operands in WGMMA. To address the decrease in computational accuracy associated with FP8 precision, we employ block quantization and incoherent processing methods.
\end{enumerate}

To empirically assess our approach, we conduct a benchmarking study of \textsc{FlashAttention-3} on the H100 SXM5 GPU across multiple parameter settings. Our results demonstrate that (1) using FP16, \textsc{FlashAttention-3} delivers a 1.5--2.0$\times$ improvement in speed over \textsc{FlashAttention-2} during the forward computation (reaching up to 740 TFLOPs/s) and achieves 1.5--1.75$\times$ acceleration for the backward computation, (2) in FP8 mode, performance approaches 1.2 PFLOPs/s, and (3) FP16 shows superior throughput and FP8 remains competitive for extensive sequence lengths\footnote{Specifically, with head dimension 64, FP8 \textsc{FlashAttention-3} is ahead; with head dimensions 128 and 256, it matches performance for non-causal masking and is behind in scenarios involving causal masking.} relative to the state-of-the-art attention implementation in NVIDIA’s cuDNN library. We further confirm that the numerical accuracy of FP16 \textsc{FlashAttention-3} aligns with that of \textsc{FlashAttention-2}, and exceeds that of conventional attention mechanisms, since intermediate computations (such as softmax scaling) retain FP32 precision. Additionally, FP8 \textsc{FlashAttention-3} utilizing block quantization and incoherent processing yields 2.6$\times$ enhanced accuracy compared to standard attention relying on per-tensor quantization, particularly in situations characterized by outlier features.

We are releasing \textsc{FlashAttention-3} under a permissive license\footnote{\textsc{FlashAttention-3} can be accessed at \texttt{https://github.com/Dao-AILab/flash-attention}} and intend to incorporate it into the PyTorch and Hugging Face ecosystems, thereby maximizing its accessibility and impact for researchers and developers.

\section{Background: Multi-Head Attention and GPU Characteristics}
\label{sec:background}

\subsection{Multi-Head Attention}
\label{subsec:multi_head_attn}

Let $\mathbf{Q}, \mathbf{K}, \mathbf{V} \in \mathbb{R}^{N \times d}$ be the query, key and value input sequences associated to a single head, where $N$ is the sequence length and $d$ is the head dimension. Then the attention output $\mathbf{O}$ is computed as:

\begin{equation*}
  \mathbf{S} = \alpha \mathbf{Q} \mathbf{K}^\top \in \mathbb{R}^{N \times N}, \quad \mathbf{P} = \mathrm{softmax}(\mathbf{S}) \in \mathbb{R}^{N \times N}, \quad \mathbf{O} = \mathbf{P}\mathbf{V} \in \mathbb{R}^{N \times d},
\end{equation*}

Here, the $\mathrm{softmax}$ function operates across each row, and it is common to use $\alpha = 1/\sqrt{d}$ for scaling. To mitigate numerical issues arising from exponentiation, $\mathrm{rowmax}(\mathbf{S})$ is subtracted from $\mathbf{S}$. In the case of multi-head attention (MHA), distinct projection matrices are employed for the queries, keys, and values of each head, enabling the simultaneous computation across multiple heads and data batches, which yields the complete output tensor.

Now let $\phi$ be a scalar loss function and let $\mathbf{d}(-) = \partial \phi / \partial (-)$ be notation for the gradient. Given the output gradient $\mathbf{dO} \in \mathbb{R}^{N \times d}$, we compute $\mathbf{dQ}$, $\mathbf{dK}$, and $\mathbf{dV}$ according to the chain rule as follows:

\begin{align*}
  \mathbf{dV} &= \mathbf{P}^\top \mathbf{dO} \in \mathbb{R}^{N \times d} \\
  \mathbf{dP} &= \mathbf{dO} \mathbf{V}^\top \in \mathbb{R}^{N \times N} \\
  \mathbf{dS} &= \mathrm{dsoftmax} (\mathbf{dP}) \in \mathbb{R}^{N \times N} \\
  \mathbf{dQ} &= \alpha \mathbf{dS} \mathbf{K} \in \mathbb{R}^{N \times d} \\
  \mathbf{dK} &= \alpha \mathbf{dS}^\top \mathbf{Q} \in \mathbb{R}^{N \times d},
\end{align*}

In this context, $\mathbf{d}s = (\mathrm{diag}(p) - p p^\top)\mathbf{d}p$ where $p = \mathrm{softmax}(s)$, treating $s$ as a vector. When applying this expression to each row, we denote it as $\mathrm{dsoftmax}(\mathbf{dP})$. For the backward pass in MHA, this operation is readily parallelized across all heads and batch instances.

\subsection{GPU hardware characteristics and execution model}
\label{subsec:hardware}

We present an overview of the GPU execution model pertinent to \textsc{FlashAttention-3}, emphasizing the NVIDIA Hopper architecture as a specific example of this framework.

\paragraph{Memory hierarchy:} On the GPU, different forms of memory are structured into a hierarchy of storage areas, where larger capacities typically correspond to lower bandwidth (\cref{tab:gpu-hierarchy})\footnote{\citet{luo2024benchmarking} specifies a shared memory bandwidth of 128 bytes per clock cycle per SM. This is then scaled by multiplying by 132 SMs and the boost clock frequency of 1830 MHz.}. At the highest level, global memory (GMEM)—or HBM—acts as off-chip DRAM, providing accessible storage to all streaming multiprocessors (SMs). When data resides in GMEM, it is automatically cached by the on-chip L2 cache to facilitate faster access. Beyond this, each SM is equipped with a small, high-speed, banked memory known as shared memory (SMEM), which must be managed explicitly by programmers. Finally, the register file represents the innermost memory level within each SM.

\paragraph{Thread hierarchy:} In the GPU programming paradigm, execution entities known as threads are systematically arranged into several hierarchical levels. Starting from the smallest unit, threads are aggregated into warps, each consisting of 32 threads. Subsequently, warps are combined into warpgroups, each made up of 4 consecutive warps. Moving up the hierarchy, threadblocks—also referred to as cooperative thread arrays (CTAs)—enclose these warpgroups. On some architectures, such as Hopper, threadblocks can further be grouped into threadblock clusters. Finally, grids represent the most extensive structure, encompassing all threadblocks assigned to a computation.

The two hierarchies exhibit significant interdependence. Threads belonging to a single CTA are executed concurrently on the same SM, while CTAs grouped within the same cluster are likewise dispatched simultaneously onto a common GPC. SMEM can be accessed directly by every thread inside a CTA, in contrast to each thread’s allocation of up to 256 private registers (RMEM) which are exclusive to that thread.

\paragraph{Asynchrony and warp-specialization:}

GPUs function as throughput-oriented processors, leveraging concurrent and asynchronous operations to mitigate delays stemming from memory access and execution. In the context of asynchronous memory transfers between GMEM and SMEM, Hopper introduces the Tensor Memory Accelerator (TMA), a specialized hardware module \cite[\S7.29]{cuda}. Additionally, in contrast to earlier designs like Ampere, Hopper's Tensor Core supports asynchrony and can retrieve its input operands directly from shared memory, facilitated by the warpgroup-wide WGMMA instruction \cite[\S9.7.14]{ptx}.

Enabling hardware support for asynchrony permits the implementation of warp-specialized kernels, wherein each warp within a CTA is exclusively designated as either a producer or a consumer. Consequently, these warps are restricted to issuing solely data movement or solely computation instructions. This arrangement generally enhances the compiler’s capacity to produce more efficient instruction schedules \citep{warp-specialization-2011}. Furthermore, Hopper introduces the capability to dynamically reassign registers among warpgroups using \verb|setmaxnreg| \cite[\S9.7.17.1]{ptx}, thereby permitting warps responsible for executing MMAs to be allocated a greater portion of RMEM relative to warps limited to issuing TMA, which requires only one thread.

\paragraph{Low-precision number formats:}
\label{sec:low-precision-gpu}
Contemporary GPUs are equipped with dedicated hardware components designed to efficiently handle computations with reduced numerical precision. As an illustration, the WGMMA operation is capable of utilizing the FP8 Tensor Cores found in Hopper architecture, enabling a single SM to achieve throughput that is double that of FP16 or BF16 operations.

To properly utilize FP8 WGMMA, it is crucial to be aware of the operand layout requirements. Consider performing GEMM for $A \times B^{\top}$, where $A$ is an $M\times K$ matrix and $B$ is an $N\times K$ matrix. An operand is designated as \emph{mn-major} if it is memory contiguous along the $M$ or $N$ axis (the outer dimension), while it is labeled \emph{k-major} if contiguous storage is along the $K$ axis (the inner dimension). For FP16 WGMMA, both mn-major and k-major layouts are permitted for SMEM inputs. In contrast, FP8 WGMMA restricts operand formatting to k-major only. Additionally, applications such as attention mechanisms—where multiple GEMM operations are fused within a single kernel—face the complication that mismatched layouts between FP32 accumulators and FP8 operands can impede chaining FP8 WGMMAs.

Regarding attention mechanisms, these layout constraints necessitate specific adjustments to how an FP8 algorithm is structured, as outlined in \cref{sec:algofp8}.

\subsection{Standard Attention and Flash Attention} Adopting the definitions from \citet{dao2022flashattention}, \textbf{standard attention} refers to the approach in which attention is implemented on the GPU by storing the intermediate matrices $\mathbf{S}$ and $\mathbf{P}$ in HBM. The central innovation of \textsc{FlashAttention} lies in utilizing a localized softmax reduction to circumvent the costly process of reading from and writing to these intermediate matrices, instead integrating the attention computation within a unified kernel. Here, the local softmax process corresponds to lines \ref{code-ws:softmax_start}-\ref{code-ws:softmax_end} within the consumer mainloop found in \cref{alg:flash3_wgmma_ws_only}, along with the application of rescaling to blocks of $\mathbf{O}$. A straightforward demonstration that this process correctly yields $\mathbf{O}$ is provided in \cite[\S 2.3.1]{dao2023flashattention2}.

\section{FlashAttention-3: Algorithm}
\label{sec:algo}

This section presents an overview of the \textsc{FlashAttention-3} algorithm. For clarity, our discussion centers on the forward pass, while the algorithmic details for the backward pass can be found in~\cref{sec:algo_ws_bwd}. We begin by outlining the method for incorporating warp-specialization alongside a circular SMEM buffer into the foundational framework provided by \textsc{FlashAttention-2}. Next, we detail the approach for leveraging the asynchrony inherent in WGMMA to construct a 2-stage pipeline that overlaps GEMM and softmax computations. Lastly, we discuss the adjustments required to support FP8, addressing layout compatibility and precision concerns through block quantization as well as incoherent processing techniques.

\subsection{Producer-Consumer asynchrony through warp-specialization and pingpong scheduling}
\label{sec:algo_ws}

\paragraph{Warp-specialization}
Similar to \textsc{FlashAttention-2}, the forward computation in \textsc{FlashAttention-3} is highly parallelizable across batch dimension, the count of heads, and the length of the query sequence. As a result, we focus on presenting the CTA-level perspective of the algorithm, where each CTA handles a tile $\mathbf{Q}_i$ from the query matrix and computes the corresponding output tile $\mathbf{O}_i$. For clarity, the initial exposition considers warp-specialization using a circular SMEM buffer, omitting the concurrent GEMM-softmax processing. Here, $d$ denotes the head size, $N$ is the overall sequence length, and a chosen query block size $B_r$ partitions $\mathbf{Q}$ into $T_r = \lceil \frac{N}{B_r} \rceil$ blocks, specifically $\mathbf{Q}_1, .., \mathbf{Q}_{T_r}$.

\begin{algorithm}[H]
    \caption{\small\label{alg:flash3_wgmma_ws_only}\textsc{FlashAttention-3} forward execution \textbf{absent} intra-consumer overlap -- CTA arrangement}
    \begin{algorithmic}[1]
\REQUIRE Inputs: $\mathbf{Q}_i \in \mathbb{R}^{B_r \times d}$, $\mathbf{K}, \mathbf{V} \in \mathbb{R}^{N \times d}$ stored in HBM; key block size $B_c$ such that $T_c = \lceil \frac{N}{B_c} \rceil$.
\STATE Set up a pipeline construct responsible for synchronizing barriers, equipped with an $s$-stage circular shared memory buffer.
\IF {executing in the producer warpgroup}
\STATE Release a fixed allocation of registers.
\STATE Transfer $\mathbf{Q}_i$ from HBM into shared memory.
\STATE Once the load is done, signal to the consumer that $\mathbf{Q}_i$ is now present.
\FOR{$0 \le j < T_c$}
    \STATE Await that the consumer has finished consuming the $(j\,\%\,s)$-indexed buffer segment.
    \STATE Trigger the transfer of $\mathbf{K}_j, \mathbf{V}_j$ from HBM into shared memory at buffer stage $(j\,\%\,s)$.
    \STATE Once completed, notify consumers that $\mathbf{K}_j, \mathbf{V}_j$ are ready for use.
\ENDFOR
\ELSE \STATE Reacquire a fixed count of registers, according to the consumer warp count.
\STATE Locally initialize on-chip accumulators: $\mathbf{O}_i = (0) \in \mathbb{R}^{B_r \times d}$, $\ell_i = (0) \in \mathbb{R}^{B_r}$, $m_i = (-\infty) \in \mathbb{R}^{B_r}$.
\STATE Block until $\mathbf{Q}_i$ has been copied to shared memory.
\FOR{$0 \le j < T_c$}
\STATE Wait until $\mathbf{K}_j$ becomes available in shared memory.
\STATE Execute $\mathbf{S}_i^{(j)} = \mathbf{Q}_i \mathbf{K}_j^T$ as an SS-GEMM operation. Synchronize at commit. \label{alg:ws_only_gemm1}
\STATE Save prior $m_i$ into $m_i^{\mathrm{old}}$; update $m_i \gets \mathrm{max}(m_i^{\mathrm{old}}, \mathrm{rowmax}(\mathbf{S}_i^{(j)}))$. \label{code-ws:softmax_start}
\STATE Evaluate $\widetilde{\mathbf{P}}_i^{(j)} = \mathrm{exp}(\mathbf{S}_i^{(j)} - m_i)$; update $\ell_i \gets \mathrm{exp}(m_i^{\mathrm{old}} - m_i) \ell_i + \mathrm{rowsum}(\widetilde{\mathbf{P}}_i^{(j)})$. \label{code-ws:softmax_end}
\STATE Wait until $\mathbf{V}_j$ becomes available in shared memory.
\STATE Update $\mathbf{O}_i \gets \mathrm{diag}(\mathrm{exp}(m_i^{\mathrm{old}} - m_i))^{-1} \mathbf{O}_i + \widetilde{\mathbf{P}}_i^{(j)} \mathbf{V}_j$ via RS-GEMM; commit and synchronize. \label{alg:ws_only_gemm2}
\STATE Mark the $(j\,\%\,s)$-indexed buffer segment as free for reuse by the producer.
\ENDFOR
\STATE Compute final $\mathbf{O}_i \gets \mathrm{diag}(\ell_i)^{-1} \mathbf{O}_i$ and $L_i \gets m_i + \log(\ell_i)$.
\STATE Persist $\mathbf{O}_i$ and $L_i$ back to HBM, writing as block $i$ of outputs $\mathbf{O}$ and $L$.
\ENDIF
\end{algorithmic}
\end{algorithm}

In our implementation of \cref{alg:flash3_wgmma_ws_only} on Hopper, memory (de)allocation is handled using \verb|setmaxnreg|, while TMA is employed to load $\mathbf{Q}_i$ as well as $\{ \mathbf{K}_j, \mathbf{V}_j \}_{0 \leq j < T_c}$. WGMMA is utilized to perform the GEMMs within the main consumer loop, with either the SS or RS prefix denoting whether the initial operand is fetched from shared memory or the register file. To understand the execution sequence of \cref{alg:flash3_wgmma_ws_only}, it is important to recognize that TMA loads operate asynchronously, allowing them to be initiated without waiting for preceding loads to complete. Additionally, during the initial $s$ passes of the producer mainloop, there will not be any wait instructions, as these iterations correspond to the period when the buffer is being populated.

\paragraph{Pingpong scheduling} By leveraging the asynchronous execution modes of WGMMA and TMA in conjunction with warp-specialization, it becomes possible to concurrently execute the softmax operation for one warpgroup while another warpgroup carries out GEMM. This is particularly appealing given the considerable performance disparity between matmul and non-matmul operations on state-of-the-art hardware accelerators. For instance, the H100 SXM5 GPU achieves up to 989 TFLOPS for FP16 matmul, yet only 3.9 TFLOPS for special function computations such as exponentials\footnote{According to the CUDA programming guide, each streaming multiprocessor (SM) is capable of executing 16 special function instructions per clock cycle. With 132 SMs operating at a clock speed of 1830 MHz, this results in an aggregate throughput of 3.9 TFLOPS for special functions.}, which are essential for softmax evaluation. Considering FP16 attention with a head dimension of 128, matmul operations account for 512 times more FLOPS than exponentials, but the latter operate at just $1/256$th of the throughput, which means exponentials might occupy up to half the computation cycles when compared to matmul. This disparity grows further in FP8, where the performance of matmul doubles but the throughput for exponential remains unchanged.

Given that the exponential operation is handled by a dedicated hardware component (the multi-function unit), optimal performance would require overlapping the exponential computation with the matrix multiplication activities performed by the Tensor Cores. To facilitate this overlap, we employ synchronization barriers (\texttt{bar.sync} instructions) so that the GEMMs (specifically, GEMM1 -- $\mathbf{P} \mathbf{V}$ from a given iteration and GEMM0 -- $\mathbf{Q} \mathbf{K}^\top$ from the subsequent iteration) within warpgroup 1 are executed prior to those of warpgroup 2. This ensures that while warpgroup 2 is engaged in its GEMMs, the softmax for warpgroup 1 is calculated concurrently. Afterwards, the execution alternates, with warpgroup 2 carrying out its softmax computation as warpgroup 1 handles the GEMMs, thereby implementing a "pingpong" style scheduling as visualized in~\cref{fig:pingpong_scheduling}. Although the real-world scheduling behavior may deviate from the idealized diagram, this technique tends to yield notable performance gains (for instance, throughput increases from 570 TFLOPS to around 620-640 TFLOPS in the FP16 forward pass when using a head dimension of 128 and sequence length of 8192).

\paragraph{Attention variants} In the cases of multi-query attention~\citep{shazeer2019fast} and grouped query attention~\citep{ainslie2023gqa}, our implementation mirrors the methodology utilized in \textsc{FlashAttention-2}. Specifically, we modify tensor indexing so that the $\mathbf{K}$ and $\mathbf{V}$ tensors are not redundantly stored in HBM.

\subsection{Intra-warpgroup overlapping GEMMs and softmax}

Within a single warpgroup, it is possible to execute certain instructions from the softmax operation simultaneously with instructions from the GEMMs. Here, we outline a specific method to achieve this overlap.

In the attention algorithm, operations within the inner loop (main loop) have sequential dependencies that impede parallelization within a single iteration. For example, (local) softmax (lines \ref{code-ws:softmax_start} to \ref{code-ws:softmax_end}) relies on the output $\mathbf{S}_i^{(j)}$ of the first GEMM, while the second GEMM takes its result $\widetilde{\mathbf{P}}_i^{(j)}$ as an operand. Indeed, the wait statements in lines \ref{alg:ws_only_gemm1} and \ref{alg:ws_only_gemm2} of \cref{alg:flash3_wgmma_ws_only} serialize the execution of softmax and GEMMs. However, we can break these dependencies by pipelining across iterations through additional buffers in registers. Pursuing this idea, we propose the following two-stage\footnote{Note that the number of stages of the overlapping scheme is bounded by, but need not equal, the number $s$ of stages in the circular SMEM buffer.} GEMM-softmax pipelining algorithm:

\begin{algorithm}[H]
    \caption{\small\label{alg:flash3_wgmma}\textsc{FlashAttention-3} consumer warpgroup forward pass}
    \begin{algorithmic}[1]
      \REQUIRE  Matrices $\mathbf{Q}_i \in \mathbb{R}^{B_r \times d}$ and $\mathbf{K}, \mathbf{V} \in \mathbb{R}^{N \times d}$ located in HBM, a specified key block size $B_c$, and $T_c = \lceil \frac{N}{B_c} \rceil$ blocks.
      \STATE Adjust allocation of registers dynamically according to the total number of consumer warps.
      \STATE Set $\mathbf{O}_i = (0) \in \mathbb{R}^{B_r \times d}$, $\ell_i = (0)$, and $m_i = (-\infty)$ within on-chip memory resources.
      \STATE Ensure that shared memory contains loaded $\mathbf{Q}_i$ and $\mathbf{K}_0$ before proceeding.
      \STATE Use WGMMA to compute $\mathbf{S}_{\mathrm{cur}} = \mathbf{Q}_i \mathbf{K}_0^T$; commit this calculation and synchronize.
      \STATE Open the $0$th stage of $\mathbf{K}$ buffer for subsequent operations.
      \STATE Determine $m_{i}$, $\tilde{\mathbf{P}}_{\mathrm{cur}}$, and $\ell_{i}$ utilizing $\mathbf{S}_{\mathrm{cur}}$, and apply a scaling step to $\mathbf{O}_i$.
      \FOR{$1 \le j < T_c - 1$}
        \STATE Confirm availability of $\mathbf{K}_j$ in shared memory.
        \label{code:mainloop_start}
        \STATE Employ WGMMA to compute $\mathbf{S}_{\mathrm{next}} = \mathbf{Q}_i \mathbf{K}_{j}^T$; initiate commit without blocking for completion. \label{code:first_wgmma}
        \STATE Verify that $\mathbf{V}_{j-1}$ is now accessible in shared memory.
        \STATE Use WGMMA to update $\mathbf{O}_{i} = \mathbf{O}_{i} + \tilde{\mathbf{P}}_{\mathrm{cur}} \mathbf{V}_{j-1}$; commit but allow nonblocking execution. \label{code:second_wgmma}
        \STATE Wait for WGMMA to finish $\mathbf{Q}_i \mathbf{K}_{j}^T$ operation.
        \STATE Calculate new values for $m_{i}$, $\tilde{\mathbf{P}}_{\mathrm{next}}$, and $\ell_{i}$ based on $\mathbf{S}_{\mathrm{next}}$. \label{code:softmax}
        \STATE After confirming WGMMA completion for $\tilde{\mathbf{P}}_{\mathrm{cur}} \mathbf{V}_{j-1}$, rescale $\mathbf{O}_i$ accordingly.
        \STATE Unlock buffer stage $(j\,\%\,s)$ for $\mathbf{K}$ and $(j-1\,\%\,s)$ for $\mathbf{V}$ to allow further access.
        \STATE Transfer contents from $\mathbf{S}_{\mathrm{next}}$ to $\mathbf{S}_{\mathrm{cur}}$ to maintain state for next iteration.
        \label{code:mainloop_end}
      \ENDFOR \STATE Await loading of $\mathbf{V}_{T_c - 1}$ into shared memory.
      \STATE Finalize output as $\mathbf{O}_{i} = \mathbf{O}_{i} + \tilde{\mathbf{P}}_{\mathrm{last}} \mathbf{V}_{T_c - 1}$ via WGMMA; commit changes and ensure completion.

\STATE Epilogue: Adjust the scale of $\mathbf{O}_{i}$ using $m_{i}$. Determine $L_{i}$ from $m_{i}$ and $\ell_{i}$.  
Store $\mathbf{O}_{i}$ and $L_{i}$ in HBM as the $i$-th segments of $\mathbf{O}$ and $L$.  
\end{algorithmic}
\end{algorithm}

\cref{alg:flash3_wgmma} substitutes for the consumer pathway found in \cref{alg:flash3_wgmma_ws_only}, thereby completing the full \textsc{FlashAttention-3} procedure for computations in FP16 format. Broadly speaking, we use WGMMA to denote asynchronous GEMM operations. In the main loop (spanning lines \ref{code:mainloop_start} through \ref{code:mainloop_end}), the second WGMMA step occurring in iteration $j$ (see line \ref{code:second_wgmma}) is executed concurrently with the softmax computations belonging to iteration $j+1$ (refer to line \ref{code:softmax}).

Although the pipelined scheme presented promises performance improvements in theory, several practical considerations must be addressed:

\paragraph{Compiler reordering} Although the pseudocode suggests an idealized sequence of operations, the NVCC compiler typically reorders instructions to enhance optimization. Such rearrangement may disturb the intended sequence of WGMMA and non-WGMMA instructions, possibly causing unanticipated results or reducing expected speedup. Examination of the generated SASS code reveals that overlapping instructions are indeed produced, as detailed in Section~\ref{sec:2-stage-sass}.

\paragraph{Register pressure} Achieving maximum throughput requires keeping register spilling to a minimum. The two-stage pipeline, however, demands extra registers to save intermediate outputs and preserve state across stages. Specifically, storing an additional $\mathbf{S}_{\mathrm{next}}$ within registers results in an increased register allocation of $B_r \times B_c \times \text{sizeof}(\text{float})$ per threadblock. This higher register usage can limit the possibility of utilizing larger threadblock sizes, which themselves rely heavily on register resources. Thus, the selection of pipeline configuration and tile size must be informed by empirical profiling to optimize trade-offs.

\paragraph{3-stage pipelining} Building upon the earlier two-stage approach, a three-stage variant is proposed, aiming to overlap the second WGMMA step with the softmax computation. While this may lead to superior utilization of the Tensor Core units, the extra pipeline stage further increases register requirements. As a result, balancing the pipeline depth against tile size becomes more challenging. Comprehensive description and analysis of this three-stage approach appear in ~\cref{sec:3-stage}.

\subsection{Low-precision with FP8}
\label{sec:algofp8}

\textbf{Efficiency: layout transformations.} Performing the forward pass of \textsc{FlashAttention-3} using FP8 precision introduces unique difficulties related to layout compatibility that do not arise with FP16.

To begin, observe that the input tensors $\mathbf{Q}$, $\mathbf{K}$, and $\mathbf{V}$ are customarily stored with contiguity along the head dimension. In contrast, to adhere to the k-major requirement of FP8 WGMMA in the second GEMM, we require $\mathbf{V}$—specifically, its tiles that are transferred into SMEM—to possess contiguity in the sequence length dimension. As the TMA load operation preserves the original contiguous axis, adjustments must be made: either (1) transpose $\mathbf{V}$ within GMEM prior to computation, or (2) perform an intra-kernel transpose on the tiles of $\mathbf{V}$ after they are placed in SMEM. For the first approach, implementation can proceed by (1a) fusing the transpose into the epilogue of a previous operation such as rotary embedding, or (1b) deploying a dedicated preprocessing transpose kernel\footnote{A highly efficient transpose kernel can reach performance close to device bandwidth~\citep{colfax_cutlass_transpose_2024}.} to switch the order of the sequence length and head axes. Nonetheless, option (1a) poses integration challenges for standardized library routines, and option (1b) incurs excessive overhead under memory-bound workloads, such as those encountered during inference.

For FP8 \textsc{FlashAttention-3}, we select approach (2) as our method of choice. To achieve the in-kernel transpose, we exploit the LDSM (\verb|ldmatrix|) and STSM (\verb|stmatrix|) instructions, which enable a group of threads within a warp to efficiently transfer data between SMEM and RMEM in contiguous 128-byte segments.%
\footnote{The PTX documentation specifies that LDSM/STSM are intended for moving $8 \times 8$ matrices with 16-bit elements \cite[\S 9.7.13.4.15-16]{ptx}. Nevertheless, it is possible to utilize LDSM/STSM for FP8 computations by packing two 8-bit entries together. However, the transpose forms of LDSM/STSM do not accommodate splitting these packed 8-bit values, thus certain register shuffles are needed between LDSM and STSM operations to complete a tile-wise transpose, the specifics of which are omitted here.} These instructions offer substantial register efficiency, making them amenable to execution within the producer warpgroup, and also provide the ability to reorder layouts during memory operations. Additionally, beginning with the second iteration, we schedule the transposition of the subsequent $\mathbf{V}$ tile so it occurs concurrently—i.e., overlapped with—the two WGMMAs processing the prior $\mathbf{V}$ and the current $\mathbf{K}$ tile.

Secondly, we note a distinction between the memory layout of the FP32 accumulator resulting from an FP8 WGMMA and that of operand A's expected arrangement when loaded into registers; this contrast does not occur in the FP16 scenario. Partial visualizations of these differing layouts can be found in \cref{fig:rmem_accum} and \cref{fig:rmem_operand}, in which register storage per thread adheres to the indicated sequence. Through the application of byte permute instructions, it becomes possible to reshape the accumulator from the initial WGMMA into a configuration compatible with subsequent WGMMA computations, aligning as well with the structure of the $\mathbf{V}$ tile generated via the in-kernel transpose. Referring to \cref{fig:rmem_accum}, the register ordering is adapted to
$$\{ \verb|d0 d1 d4 d5 d2 d3 d6 d7| \},$$
with this specific register rearrangement performed repetitively for each group of 8 bytes. Conceptually, this adjustment permutes the columns within the $\mathbf{P}$ tile (so columns $0189$ become the initial columns). To produce the desired output tile via WGMMA, the in-kernel transpose can be orchestrated to yield a corresponding permutation of the rows in the $\mathbf{V}$ tile.\footnote{By employing the in-kernel transpose in this way, we gain the flexibility to avoid utilizing shuffle instructions for redistributing register ownership among threads, in contrast to our prior approach discussed in~\citep{colfax_fp8_flashattention_2024}.}

\textbf{Accuracy: block quantization and incoherent processing.} The FP8 (e4m3) representation allocates only 3 bits to the mantissa and 4 bits for the exponent, which inherently leads to increased numerical inaccuracy compared to FP16 or BF16 formats. Additionally, large-scale models are prone to the presence of outlier values~\citep{dettmers2208llm,
  sun2024massive} whose magnitudes far exceed typical values, complicating the quantization process. A common approach to mitigate this is per-tensor scaling~\citep{micikevicius2022fp8}, where a single scaling parameter is maintained for each tensor (for example, one each for $\mathbf{Q}$, $\mathbf{K}$, and $\mathbf{V}$). In this work, we incorporate two methods to address the elevated numerical error associated with FP8 during attention computations:

\begin{enumerate}

\item \textbf{Block quantization}: in this method, a single scalar value is retained for each block, such that for tensors $\mathbf{Q}$, $\mathbf{K}$, and $\mathbf{V}$, we partition them into blocks of dimensions $B_r \times d$ or $B_c \times d$ and perform quantization independently for each block. This quantization can be seamlessly integrated with operations immediately preceding the attention step (such as rotary embedding), without incurring additional latency, given that rotary embedding is limited by memory bandwidth. As the \textsc{FlashAttention-3} algorithm inherently processes data in blocks, rescaling each block of $\mathbf{S}$ to accommodate this block-wise quantization incurs no extra computational expense.

\item \textbf{Incoherent processing}: to mitigate the effect of anomalous values, both $\mathbf{Q}$ and $\mathbf{K}$ are multiplied by a randomly chosen orthogonal matrix $\mathbf{M}$ prior to FP8 quantization. Because the orthogonality of $\mathbf{M}$ ensures $\mathbf{M} \mathbf{M}^\top = I$, we have $(\mathbf{Q} \mathbf{M}) (\mathbf{K} \mathbf{M})^\top = \mathbf{Q} \mathbf{K}^\top$, meaning the output of the attention mechanism remains unaffected by this transformation. This approach helps distribute outliers across entries, as each element in $\mathbf{Q} \mathbf{M}$ or $\mathbf{K} \mathbf{M}$ becomes a random linear combination of original values, thereby decreasing quantization error. Following the methodologies of \citet{chee2024quip} and~\citet{tseng2024quip}, $\mathbf{M}$ is constructed as a product of random diagonal matrices containing $\pm 1$ and a Hadamard matrix, facilitating multiplication in $O(d \log d)$ rather than $O(d^2)$, and this transformation can also be combined with rotary embedding without additional computational overhead.
\end{enumerate}

Empirically, these two strategies can decrease numerical error by up to a factor of 2.6, as demonstrated in \cref{sec:numerical_error}.

\section{Empirical Validation}
\label{sec:exp}

Our implementation of \textsc{FlashAttention-3} leverages WGMMA and TMA abstractions provided by CUTLASS~\citep{Thakkar_CUTLASS_2023} primitives, which form the basis for our assessments of both performance and correctness.

\begin{itemize}

\item \textbf{Benchmarking attention.} We evaluate the performance of \textsc{FlashAttention-3} by recording its runtime for various sequence lengths and performing a comparative analysis against several baselines: the standard PyTorch attention implementation, \textsc{FlashAttention-2}, \textsc{FlashAttention-2} implemented with Triton (leveraging H100-specific instructions), and a cuDNN-optimized vendor version of \textsc{FlashAttention-2} for H100 GPUs. Our results show that \textsc{FlashAttention-3} achieves speedups of up to 2.0$\times$ compared to \textsc{FlashAttention-2} and is 1.5$\times$ faster than the Triton implementation. Notably, \textsc{FlashAttention-3} attains peak throughput of 740 TFLOPs/s, corresponding to 75\% of the H100 GPU's theoretical maximum.

\item \textbf{Ablation study.} Through detailed ablations, we demonstrate that the integration of warp-specialization techniques and pipelining of GEMM-softmax operations are major contributors to the performance enhancements seen in \textsc{FlashAttention-3}.

\item \textbf{Accuracy of FP8 attention.} Our experiments indicate that employing block quantization along with incoherent processing substantially reduces the numerical error in FP8 \textsc{FlashAttention-3}, yielding a 2.6$\times$ decrease in error.
\end{itemize}

\subsection{Benchmarking Attention}
\label{subsec:benchmark_attn}

We evaluate the execution time of various attention mechanisms on an H100 80GB SXM5 GPU, using FP16 data and considering multiple configurations (with or without a causal mask, and head dimensions set to 64 or 128). The benchmarking outcomes are displayed in~\cref{fig:benchmark_attn_fwd} and~\cref{fig:benchmark_attn_bwd}. The data indicates that \textsc{FlashAttention-3} achieves approximately 1.5-2.0$\times$ speedup over \textsc{FlashAttention-2} during the forward operation and is 1.5-1.75$\times$ faster during the backward operation. When measured against a typical attention implementation, \textsc{FlashAttention-3} demonstrates acceleration of up to 3-16$\times$. Particularly for sequences of moderate to large length (starting from 1k tokens), \textsc{FlashAttention-3} even outperforms the proprietary vendor library cuDNN—which is specifically optimized for H100 hardware.

\paragraph{Benchmark settings:} We vary the sequence length as 512, 1k, ..., 16k, and set batch size so that the total number of tokens is 16k. We set the hidden dimension to 2048, and head dimension to be either 64, 128, or 256 (i.e., 32 heads, 16 heads, or 8 heads). To calculate the FLOPs of the forward pass, we use:

\begin{equation*}
  4 \cdot \text{seqlen}^2 \cdot \text{head dimension} \cdot \text{number of heads}.
\end{equation*}

When employing causal masking, we adjust this figure by halving it, reflecting that roughly 50% of the elements are processed. To estimate the FLOPs for the backward pass, we scale the forward pass FLOPs by a factor of 2.5. This multiplier arises because the forward computation involves 2 matrix multiplications, whereas the backward computation (with recomputation) encompasses 5 matrix multiplications.

Additionally, we evaluate the forward pass runtime using FP8 under comparable conditions. The outcomes for headdim 256 are presented in ~\cref{fig:benchmark_attn_fp8}, with comprehensive results provided in \cref{sec:benchmark_attn_fp8_full}.

\subsection{Ablation Study: 2-Stage Pipelining Experiments}
\label{sec:2-stage-pipelining-experiments}

We conduct an ablation study on the 2-stage WGMMA-softmax pipelining and the warp-specialization technique within the context of non-causal FP16 \textsc{FlashAttention-3}, holding the following parameters constant: $\{\text{batch}, \text{seqlen}, \text{nheads}, \text{hdim}\} = \{ 4, 8448, 16, 128\}$. 
The outcomes reported in~\cref{table:ablation_pipelining} demonstrate that the algorithmic enhancements we introduce—asynchronous execution paired with warp-specialization as well as concurrent execution of GEMM and softmax—substantially boost performance, increasing throughput from 570 to 661 TFLOPs.

\subsection{Numerical Error Validation}
\label{sec:numerical_error}

Due to ongoing concerns regarding the numerical inaccuracies of \textsc{FlashAttention}~\citep{golden2024flash}, we evaluate \textsc{FlashAttention-2}, \textsc{FlashAttention-3}, and a conventional attention implementation by benchmarking each against a reference computation performed in FP64 precision. In order to mimic the occurrence of outlier activations and feature values commonly seen in LLMs~\citep{dettmers2208llm, sun2024massive}, we sample the elements of $\mathbf{Q}, \mathbf{K}, \mathbf{V}$ according to the following probability distribution:

\begin{equation*}
  \mathcal{N}(0, 1) + \mathcal{N}(0, 100) \cdot \mathrm{Bernoulli}(0.001).
\end{equation*}

Specifically, each element is sampled from a normal distribution with a mean of zero and a standard deviation of 1; however, for 0.1\% of the elements, we add an additional independent noise component drawn from a normal distribution with a standard deviation of 10. The root mean squared error (RMSE) under these conditions is presented in~\cref{table:numerical_error}. For computations in FP16, both \textsc{FlashAttention-2} and \textsc{FlashAttention-3} yield a RMSE that is 1.7$\times$ lower than that of the standard method, as they store intermediate values (such as softmax outputs) in FP32. For FP8, the baseline attention approach employs per-tensor scaling, with accumulation during matrix multiplication performed in FP32 and intermediate softmax results stored in FP16. Due to the incorporation of block quantization and incoherent processing techniques, \textsc{FlashAttention-3} operating in FP8 attains 2.6$\times$ higher accuracy relative to this baseline.

\section{Dicussion, Limitations, Conclusion}
\label{sec:discussion}

By employing \textsc{FlashAttention-3}, we have shown that adopting novel programming strategies and leveraging hardware advancements, including asynchronous computation and reduced-precision formats, significantly enhance both the performance and precision of attention mechanisms. Compared to \textsc{FlashAttention-2}, our implementation accelerates attention computation by 1.5–2.0$\times$, and lowers FP8 numerical error by a factor of 2.6 relative to conventional per-tensor quantization. Our current approach has a few shortcomings, which we aim to resolve in subsequent work, such as further optimizing for LLM inference scenarios, incorporating a persistent kernel into the FP8 kernel,\footnote{Within our benchmarks, FP16 \textsc{FlashAttention-3} utilizes a persistent kernel and implements load balancing, whereas FP8 \textsc{FlashAttention-3} lacks these features. This is in part responsible for the lower performance of FP8 \textsc{FlashAttention-3} on tasks with short sequence lengths and causal masking, especially compared to FP8 cuDNN kernels.} and thoroughly evaluating the consequences of low-precision attention in large-scale model training. Although our experiments primarily target Hopper GPUs, the underlying methodologies should be transferable to a wide range of hardware accelerators. Ultimately, we anticipate that refining attention as a fundamental operation will broaden the horizon for efficient, long-context applications.

\end{document}